%% Elevark fra mattelab.no – laget automatisk av bygg.py. Rediger gjerne til eget bruk.
%% Kompiler med lualatex.
\documentclass[11pt,a4paper]{article}
\usepackage[norsk]{babel}
\usepackage[a4paper,margin=2.2cm,top=2cm,bottom=2.4cm]{geometry}
\usepackage{fontspec}
\usepackage{unicode-math}
\setmainfont{TeX Gyre Pagella}
\setsansfont{TeX Gyre Heros}[Scale=MatchLowercase]
\setmathfont{TeX Gyre Pagella Math}
\usepackage{xcolor}
\definecolor{blekk}{HTML}{16213A}
\definecolor{dempet}{HTML}{566077}
\definecolor{rute}{HTML}{CFC9B8}
\definecolor{m1}{HTML}{2F4AB0}
\definecolor{m2}{HTML}{99266D}
\usepackage[most]{tcolorbox}
\usepackage{qrcode}
\usepackage{needspace}
\usepackage{fancyhdr}
\usepackage[hidelinks]{hyperref}
\pdfvariable suppressoptionalinfo \numexpr 1+2+4+8+16+32+64+128+256+512\relax
\setlength{\parindent}{0pt}
\setlength{\parskip}{0.4em}
\pagestyle{fancy}
\fancyhf{}
\renewcommand{\headrulewidth}{0pt}
\fancyfoot[L]{\sffamily\footnotesize\color{dempet}mattelab.no · Polynomfunksjoner · elevark uten fasit}
\fancyfoot[R]{\sffamily\footnotesize\color{dempet}Side \thepage}

\newcommand{\fase}[2]{\tcbox[on line, colframe=#2, colback=white, coltext=#2, boxrule=0.6pt, arc=2pt,
  boxsep=0pt, left=3pt, right=3pt, top=1.5pt, bottom=1.5pt, fontupper=\sffamily\bfseries\footnotesize]{#1}}
\newcommand{\steg}[4]{%
  \par\noindent{\color{rute}\rule{\linewidth}{0.4pt}}\par\nopagebreak
  \noindent\begin{minipage}[t]{2.1cm}\vspace{0pt}\fase{#1}{#2}\end{minipage}%
  \begin{minipage}[t]{\dimexpr\linewidth-2.1cm\relax}\vspace{0pt}#3\end{minipage}\par
  \vspace{#4}}
\newcommand{\del}[4]{%
  \par\noindent{\color{rute}\rule{\linewidth}{0.4pt}}\par\nopagebreak
  \noindent\begin{minipage}[t]{1.1cm}\vspace{0pt}\fase{#1}{#2}\end{minipage}%
  \begin{minipage}[t]{\dimexpr\linewidth-1.1cm\relax}\vspace{0pt}#3\end{minipage}\par
  \vspace{#4}}
\newcommand{\oppgave}[2]{\par\Needspace{12\baselineskip}\vspace{0.6em}%
  {\fontsize{14}{17}\selectfont\bfseries{\color{m1}#1}\enspace #2}\par\vspace{0.2em}}

\begin{document}
\begin{minipage}[t]{0.72\linewidth}\vspace{0pt}
{\sffamily\small\color{dempet}Elevark · 1T · mattelab.no}\par\vspace{0.3em}
{\fontsize{24}{28}\selectfont\bfseries Polynomfunksjoner}\par\vspace{1em}
{\sffamily\small Navn: \rule{4.6cm}{0.4pt}\quad Klasse: \rule{1.5cm}{0.4pt}\quad Dato: \rule{1.9cm}{0.4pt}}
\end{minipage}\hfill
\begin{minipage}[t]{0.22\linewidth}\vspace{0pt}\raggedleft
\qrcode[height=2.3cm]{https://mattelab.no/utforskere/polynomfunksjoner.html}\par\vspace{0.2em}
{\sffamily\scriptsize\color{dempet}Utforskeren}
\end{minipage}

\vspace{1em}
{\sffamily\small\color{dempet}Utforskeren ligger på \url{https://mattelab.no/utforskere/polynomfunksjoner.html}}\par
Hver oppgave har tre steg. \textbf{Gjett} først, uten å bruke utforskeren. \textbf{Test} så i utforskeren. \textbf{Forklar} til slutt med egne ord hvorfor det blir slik.

\oppgave{1}{Hvordan ender grafen?}
\steg{Gjett}{m1}{Se på $f(x) = a(x + 2)(x - 1)(x - 3)$. Hva tror du skjer med $f(x)$ når $x$ blir veldig stor? Når $x$ blir veldig negativ? Skisser grafen for $a = 1$ og for $a = -1$.}{2.2cm}
\steg{Test}{m2}{Sett grad $3$ med $a = 1$. Dra $a$ over til $-1$. Bytt så til grad $4$ med nullpunktet $r_4 = 4$ i tillegg, og gjør det samme.}{2.6cm}
\steg{Forklar}{blekk}{Hvorfor er det bare graden og fortegnet til $a$ som avgjør hvordan grafen ender?}{3.4cm}

\oppgave{2}{Fra produkt til sum}
\steg{Gjett}{m1}{Hvor skjærer grafen til $f(x) = (x + 2)(x - 1)(x - 3)$ y-aksen? Prøv å svare uten å gange ut.}{2.2cm}
\steg{Test}{m2}{Sett opp funksjonen. Les av det utvidede uttrykket og skjæringen med y-aksen. Gang ut for hånd og kontroller.}{2.6cm}
\steg{Forklar}{blekk}{Vis at konstantleddet er $a \cdot (-r_1)(-r_2)(-r_3)$ for en tredjegradsfunksjon. Hva betyr det når ett av nullpunktene er $0$?}{3.4cm}

\oppgave{3}{Et dobbelt nullpunkt}
\steg{Gjett}{m1}{Se på $f(x) = (x - 1)^2(x + 2)$. Hva tror du grafen gjør i $x = 1$? Krysser den x-aksen?}{2.2cm}
\steg{Test}{m2}{Sett $r_1 = r_2 = 1$ og $r_3 = -2$. Se på fortegnsskjemaet. Dra $r_2$ litt bort fra $1$ og tilbake. Hva skjer med grafen nær $x = 1$?}{2.6cm}
\steg{Forklar}{blekk}{Hvorfor skifter ikke $f$ fortegn i $x = 1$? Hvorfor må $(1, 0)$ være et bunnpunkt?}{3.4cm}

\oppgave{4}{Finn uttrykket}
\steg{Gjett}{m1}{En tredjegradsfunksjon har nullpunktene $-1$, $2$ og $4$, og grafen skjærer y-aksen i $(0, 4)$. Hva er $a$?}{2.2cm}
\steg{Test}{m2}{Sett grad $3$ og nullpunktene $-1$, $2$ og $4$. Dra $a$ til grafen går gjennom $(0, 4)$. Kontroller med svaret.}{2.6cm}
\steg{Forklar}{blekk}{Sett opp en likning for $a$ og løs den. Skriv funksjonen både på faktorisert og utvidet form.}{3.4cm}

\oppgave{5}{Hvor mange topp- og bunnpunkter?}
\steg{Gjett}{m1}{Hvor mange topp- og bunnpunkter kan en tredjegradsfunksjon ha? Og en fjerdegradsfunksjon? Kan en tredjegradsfunksjon ha ingen?}{2.2cm}
\steg{Test}{m2}{Se på $f(x) = x(x - 3)(x + 3)$, $g(x) = (x - 1)^3$ og $h(x) = (x + 2)(x + 1)(x - 1)(x - 2)$. Skriv ned topp- og bunnpunktene.}{2.6cm}
\steg{Forklar}{blekk}{Bruk den deriverte til å forklare hvorfor en funksjon av grad $n$ har høyst $n - 1$ topp- og bunnpunkter.}{3.4cm}

\end{document}
